% PLEASE USE THIS FILE AS A TEMPLATE FOR THE PUBLICATION
% Check file IOS-Book-Article.tex
%

\documentclass{IOS-Book-Article}     %[seceqn,secfloat,secthm]
% \usepackage{mathptmx}
% \usepackage[T1]{fontenc}
% \usepackage{times}
%
%%%%%%%%%%% Put your definitions here

\usepackage{listings}
\usepackage{xcolor}
\usepackage{amsmath}
\usepackage{amssymb}
\usepackage{algorithm}
\usepackage{algpseudocode}
\usepackage{adjustbox}
\usepackage{array}

% TODO:
%
% Existing synthetic benchmarks:
% Je kan zeggen dat ze niet gemaakt zijn met neuro-symbolic reasoners in mind, er inherente biases zijn naar single-hop reasoning, ze hardcoded zijn voor een enkele ontologie en geen negative sampling aanbieden
%
% 

% "Standard data generators output saturated graphs, natively lacking the 'Base vs. Inferred' split required to train neural reasoners. To create a forward-chaining baseline, we had to artificially mask the data or run inference, whereas \textsc{Synthology} inherently generates this required split by design.

% State-of-the-art Datalog engines like RDFox are peerless in their ability to execute forward-chaining materialization. However, using them to synthesize training data for neural models requires randomly generating massive ABoxes, which results in sparse multi-hop paths, a high ratio of dead-end noise, and an inability to natively generate proof-based negative samples. \textsc{Synthology} is not designed to replace RDFox as a reasoner; it is designed to replace random forward-chaining pipelines. By utilizing backward-chaining, \textsc{Synthology} acts as a precision data-engineering tool, purposefully crafting the exact multi-hop structures and hard negative samples that neuro-symbolic models require to learn complex ontologies.

\newcommand{\notdone}{\textcolor{red}{\textbf{TODO}}}
\newcommand{\todo}[1]{\textcolor{red}{\textbf{TODO:} #1}}

\renewcommand{\P}{\mathbb{P}}

% ! Always give a \label where possible and use \ref for cross-referencing.

% Code listing style for Turtle/Ontologies
\lstdefinelanguage{turtle}{
  keywords={@prefix, a, owl:Class, owl:ObjectProperty, rdfs:subClassOf, rdfs:domain, rdfs:range},
  keywordstyle=\color{blue}\bfseries,
  identifierstyle=\color{black},
  sensitive=false,
  comment=[l]{\#},
  commentstyle=\color{gray}\ttfamily,
  stringstyle=\color{red}\ttfamily,
  morestring=[b]",
}
\lstset{
    basicstyle=\ttfamily\small,
    breaklines=true,
    frame=single,
    numbers=left,
    numberstyle=\tiny\color{gray}
}

% Commands
\newcommand{\R}{\mathbb{R}}
\newcommand{\gen}[1]{\langle #1 \rangle}
\newcommand{\KB}{\mathrm{KB}}


%%%%%%%%%%% End of definitions
\begin{document}
\pagestyle{plain}
\begin{frontmatter}          % The preamble begins here.
	%
	%\pretitle{}
	\title{Synthology: Ontology-Based Data Generation for Link Prediction Models}
	\runningtitle{Synthology: Ontology-Based Data Generation for Link Prediction Models}
	%\subtitle{}

	% For one author:
	\author[A]{\fnms{Vincent} \snm{Van Schependom}}%
	\thanks{Corresponding Author: },
	\author[B]{\fnms{Cas} \snm{Proost}} and
	\author[C]{\fnms{Pieter} \snm{Bonte}}
	\address{Department of Computer Science, KU Leuven campus Kulak Kortrijk}

	% \address[A]{Promoter}
	% \address[B]{Supervising PhD}

	\runningauthor{V. Van Schependom et al.}

	\begin{abstract}
		Knowledge Graphs (KGs) and ontologies provide the foundational structure for representing and reasoning over domain-specific knowledge in the Semantic Web. Traditional rule-based symbolic reasoners deliver perfect logical soundness and completeness, yet they are computationally expensive on large-scale graphs and brittle in the presence of noise or incompleteness, both of which are problems that occur in real-world data. Neuro-symbolic AI has therefore emerged to combine the scalability and noise robustness of neural networks with the logical consistency of symbolic reasoning.

		A persistent bottleneck, however, is the scarcity of high-quality training data for neuro-symbolic link prediction. Unlike purely symbolic reasoners, which can be applied out-of-the-box to any ontology, neuro-symbolic models must first be trained on data that faithfully follows the specific logical rules of the target ontology. Existing real-world datasets are noisy and incomplete, while existing synthetic data generators lack the complex multi-hop inference chains and challenging negative samples required to unlock genuine logical capabilities.

		This paper introduces \textsc{Synthology}, a novel, ontology-agnostic data generator that solves the training-data bottleneck. By leveraging the rule structure of \textit{any} OWL2 RL ontology in a backward-chaining manner, \textsc{Synthology} purposefully engineers complete multi-hop proof paths together with high-quality, proof-based hard negatives. We evaluate the generator by training a state-of-the-art neuro-symbolic model on its produced data and demonstrate substantial gains over models trained with conventional forward-chaining baselines.
	\end{abstract}

	\begin{keyword}
		Link Prediction, Neuro-Symbolic AI, Ontology-based Data Generation, Backward-chaining, Negative Sampling
	\end{keyword}

\end{frontmatter}

\section{Introduction}

Knowledge Graphs (KGs) and ontologies form the backbone of structured knowledge representation on the Semantic Web. While traditional symbolic reasoners (\todo{reference and cite}) can enforce logical rules with perfect soundness and completeness, they become prohibitively expensive on massive, often incomplete real-world graphs and are highly sensitive to noise. Neuro-symbolic AI addresses this tension by combining the scalability and robustness of Machine Learning with the logical precision of symbolic reasoning. These models learn to perform ``approximate reasoning,'' predicting which additional facts are logically entailed by a given ontology.

To train neuro-symbolic models effectively, large volumes of high-quality data are required. However, real-world datasets (e.g., DBpedia, Freebase \todo{cite}) are typically noisy or incomplete, making it hard to disentangle genuine logical learning from statistical shortcuts. This creates a persistent data bottleneck: without training data that specifically targets and enforces multi-hop deductive reasoning, advanced neuro-symbolic models cannot fully demonstrate or develop their capacity for logical entailment.

Currently, researchers rely on synthetic data generators to fill this gap. Yet, existing approaches -- such as unguided forward-chaining materialization or Answer Set Programming (ASP) -- often produce structurally shallow graphs. They lack the complex multi-hop inference chains and challenging negative samples required to unlock genuine logical capabilities, leaving models under-prepared for the complex inference tasks they are meant to solve.

To address this bottleneck, we present \textsc{Synthology}, an ontology-agnostic synthetic data generator. \textsc{Synthology} parses the rule structure of \textit{any} given ontology written in the OWL2 RL profile and employs a backward-chaining architecture. Instead of generating random base facts, it purposefully engineers training data that inherently requires multi-hop logical deduction to be correctly predicted. Crucially, it also generates qualitative `hard' negative samples based on formal proof trees. In this way, neuro-symbolic models can be trained and evaluated on data that truly tests their reasoning capabilities.

The remainder of this paper is structured as follows. Section \ref{sec:background} provides the necessary background on KGs and neuro-symbolic reasoning. Section \ref{sec:related_work} discusses related work in knowledge graph embeddings and existing data generation methods. Section \ref{ch:methodology} details the design and architecture of \textsc{Synthology}. Finally, Sections \ref{ch:experiments} and 5 evaluate the multi-hop reasoning capabilities, negative sampling quality, and scalability of our approach by training the Recursive Reasoning Network (RRN) \cite{hohenecker2020ontology} on data generated from a range of OWL2 RL ontologies.

% =========================
% CHAPTER 2: BACKGROUND
% =========================
\section{Background}
\label{sec:background}

\subsection{Knowledge Graphs and Ontologies}
A Knowledge Graph (KG) is a directed labeled graph $\mathcal{G} = (\mathcal{V}, \mathcal{E})$ where $\mathcal{V}$ is the set of vertices (entities) and $\mathcal{E}$ is the set of edges (relations). In the context of the Semantic Web, a KG typically contains the concrete assertions or facts about specific entities, referred to as the ABox.

To provide structural semantics and enable reasoning, the KG is paired with an ontology (the TBox). The ontology defines the schema: the classes, properties, and logical rules that govern the relationships between the entities. Together, the factual KG and the ontology form a comprehensive Knowledge Base (KB). For example, the triple $\gen{\text{Alice}, \text{hasParent}, \text{Bob}}$ is a fact in the KG (ABox), while an ontological rule like
$$\text{hasParent}(X, Y) \land \text{hasParent}(Y, Z) \implies \text{hasGrandparent}(X, Z)$$
is part of the ontology (TBox).

The Web Ontology Language (OWL) is a widely used standard for representing these ontologies. Its OWL2 RL profile balances expressiveness with computational efficiency, making it highly suitable for rule-based reasoning \cite{motik2009owl}. Consequently, OWL2 RL serves as the target ontology format for our data generator.

In the context of this paper, we refer to the input facts present in the KG as \textit{base} facts, and the logically entailed facts as \textit{inferred} facts.

\subsection{Neuro-Symbolic Reasoning}
Neuro-symbolic AI merges Machine Learning (ML) with symbolic logic to perform approximate ontology reasoning. An example is the \textit{Recursive Reasoning Network} (RRN), proposed by \cite{hohenecker2020ontology}, which we utilize as our primary evaluation model, given it's explicit design and promising performance for multi-hop reasoning tasks.

The RRN is trained to predict $\P(\gen{s,p,o} = 1)$, the likelihood of a triple being logically entailed by the provided base facts. It employs an iterative message-passing mechanism, similar to Graph Neural Networks (GNNs), where entity embeddings are recursively updated based on relations between and class memberships of neighboring entities. This allows the model to capture complex, multi-hop logical dependencies, provided the training data is robust enough to teach these patterns.

\subsection{Negative Sampling}
A critical component of training both standard embedding models and advanced networks like the RRN is negative sampling. Since KGs contain only positive facts, models must learn to distinguish true triples from false ones by training against \textit{corrupted} negatives (e.g., $\gen{s, p, o'}$ where $o'$ is random).

However, simple random negatives often fail to challenge a model's reasoning capabilities, as they can be easily identified as false based on shallow heuristics. This necessitates \textit{hard} negatives: structurally similar, plausible facts that are not logically entailed. In neuro-symbolic models utilizing the (Local) Closed World Assumption ([L]CWA), like the RRN, unentailed facts are treated as strictly \textit{false}, so can use these near-misses as challenging negatives.

% =========================
% CHAPTER 3: RELATED WORK
% =========================
\section{Related Work}
\label{sec:related_work}

\subsection{Knowledge Graph Embeddings (KGEs)}
Standard Knowledge Graph Embeddings (KGEs) map entities and relations to low-dimensional vector spaces to predict missing links. Geometric models, such as TransE \cite{transe}, model relationships as translations ($\vec{s} + \vec{p} \approx \vec{o}$). While effective for simple relations, TransE struggles with complex topologies like $1$-to-$N$ or symmetric relations. Tensor factorization models, like DistMult \cite{distmult} and ComplEx, use tensor products to better capture these interactions.

However, standard KGEs generally excel only at capturing simple, single-hop patterns through vector-space geometry. Most standard KGE benchmarks (e.g., FB15k-237) predominantly contain inferrable facts requiring only 1 or 2 \textit{hops} (chained relations) to infer, which can be solved via basic pattern matching. Because these standard KGEs are inherently shallow, they struggle with the deep, multi-hop logical chains that arise in expressive OWL2 RL ontologies, motivating the shift toward neuro-symbolic architectures like the RRN.

\subsection{Existing Data Generation Methods}
To train models for logical reasoning, the community currently relies on two primary data generation paradigms: forward-chaining materialization based on random base facts and Answer Set Programming (ASP).

In the first of these approaches, generators create a randomized set of base facts (ABox) constrained only by basic domain and range restrictions. Semantic reasoners -- like such as RDFox, Apache Jena, and Kolibrie -- then execute the ontology rules forward to materialize all entailed facts. Because the initial base facts are generated entirely at random, there is no guarantee that deep, multi-hop logical chains will form. This unguided approach typically results in structurally shallow graphs with a high ratio of dead-end noise. Furthermore, even if deep, multi-hop chains are generated by chance, they are vastly outnumbered by simple, single-hop inferences. This severe class imbalance creates a training bias, causing models to overfit on shallow topological heuristics rather than effectively learning to navigate and resolve the rare multi-hop chains.

Alternatively, ASP solvers (such as Clingo or DLV) can generate logically consistent datasets by computing stable models from logic programs. While ASP ensures strict logical validity, it acts as a brute-force model generator. It offers limited control over the specific depth of the reasoning trees and lacks the native ability to purposefully engineer the targeted, high-quality negative samples required to robustly train neuro-symbolic classifiers.

% =========================
% METHODOLOGY
% =========================
% =========================
% METHODOLOGY
% =========================
\section{Methodology: Targeted Ontology-Based Data Generation}
\label{ch:methodology}

To resolve the shortage of high-quality, multi-hop reasoning data, we introduce \textsc{Synthology}, an ontology-agnostic synthetic data generator.
%Traditional data generation pipelines generally rely on the random instantiation of base facts, subsequently employing standard (forward-chaining) reasoners merely to discover and materialize the corresponding training labels. 

In contrast to standard approaches mentioned above, the core novelty of \textsc{Synthology} lies in its targeted, ontology-guided generation of base facts. It's capable to dynamically generate an ABox driven by logical necessity rather than statistical randomness.

\subsection{Architecture Overview}
\textsc{Synthology} accepts any valid OWL2 RL ontology. The system architecture consists of three primary phases: parsing, generation, and negative sampling.

Firstly, during the parsing phase, the engine extracts the TBox constraints, separating class hierarchies (e.g., \texttt{rdfs:subClassOf}), property constraints (e.g., \texttt{rdfs:domain}, \texttt{rdfs:range}), and complex OWL rules (e.g., \texttt{owl:inverseOf}, \texttt{owl:TransitiveProperty}).

Next, rather than populating an ABox with randomized base facts and hoping complex paths emerge, the system adopts a \textit{goal-driven} generation phase. It first selects a to-be-inferred fact (the desired target label), which may require a complex deduction path. It then recursively traces the TBox rules backward to synthesize the required premises. In doing so, it creates an \textit{inferred fact} (the root of the proof tree) by specifically engineering the necessary \textit{base facts} (the leaves of the proof tree). This ensures the generated datasets contain a complete logical path to support the target fact as a derivable conclusion.

Finally, in the negative sampling phase, the proof trees generated in the previous step are systematically corrupted to create hard negative examples for robust classification training. As alternative negative sampling strategies, we also implement pure random corruption and constrained random corruption (where the corrupted entity must still satisfy domain and range constraints).

\subsection{Generation Logic \& Its Customizability}

\notdone

\subsection{Proof-Based Negative Sampling}

% Good, but can be more detailed about which part of a fact is chosen to get corrupted, how many facts are corrupted, it now feels like a high-level explanation rather than the actual in-depth methodology. Also maybe add example with image of corrupted proof tree

% heeft misschien een iets meer fomrele behandeling van de corruption nodig, bijvoorbeeld of je door corruptie van een base fact effectief de base fact vervangt door de corrupted fact, of naast de ooriginele base fact nu een corrupted base fact toevoegt, of dat corruptie van een proof tree inhoudt dat je de hele tree dupliceert met andere individuals behalve dat je een base fact corrupt. Ik denk dat je het dus best behandelt door te sepcifieren hoe je een proof tree formeel voorstelt, hoe je een instantiation van een proof tree voorstelt, i.e. een set base facts en een set inferred facts, en dan wat er met die instantiation gebeurt als je ze corrupt

Random corruption involves randomly replacing the head or tail of a true triple with a random entity from the vocabulary. For neural reasoners, these trivial negatives fail to teach the model complex logical boundaries.

To address this, \textsc{Synthology} introduces \textit{Proof-Based Corruption}. By storing the exact derivation tree of a generated target fact, we can introduce a controlled logical flaw. If a 3-hop target relies on the chain $A \xrightarrow{R_1} B \xrightarrow{R_2} C \xrightarrow{R_3} D$, \textsc{Synthology} will corrupt a leaf node in the proof tree (i.e., a targeted base fact), invalidating the entire proof. The resulting fact $\gen{A, R_\text{target}, D}$ is then labeled as an unknown (unentailed) target. Because this `near-miss' shares significant localized topology with the true facts, it forces the RRN to evaluate strict logical entailment rather than relying on statistical graph proximity and simple heuristics.

In case the corruption of a leaf node results in a fact that is still logically entailed (e.g., due to multiple parallel proof paths existing in the generated graph), the system will fall back to constrained random corruption to cleanly satisfy the negative sampling ratio $\rho_{\pm}$ set in the configuration.

% =========================
% CHAPTER 4: EXPERIMENTS
% =========================
\section{Experimental setup}
\label{ch:experiments}

To validate \textsc{Synthology}, we evaluate the quality of the generated data by training the RRN \cite{hohenecker2020ontology} and measuring its ability to predict entailed links (inferred facts). We compare the performance of models trained on \textsc{Synthology} data against datasets generated by industry-standard semantic reasoners.

\subsection{The Baseline: Forward-Chaining Materialization}

As a baseline, we simulate the standard methodology for generating synthetic logic data: unguided ABox generation paired with forward-chaining materialization, as previously described.

We employ Apache Jena and Nemo as reasoning engines. We generate a set of random base facts $\mathcal{G}_\text{base}$ strictly adhering to the TBox constraints.

% Mention how this is generated eg. by selecting a random subject, then  predicate, then object, or just by generating the set of all base facts and picking some random facts. Also make sure to keep these raw datasets and mention they are available

We then load $\mathcal{G}_\text{base}$ alongside the ontology into the reasoner, executing full materialization to derive all inferred facts $\mathcal{G}_\text{inferred}$. The resulting inferred facts act as the positive targets,

To generate negative samples, we employ random corruption with a $\rho_{\pm} = 1:1$ positive-to-negative ratio. For both the base and inferred datasets, we corrupt the subject, predicate (\todo{echt?}), or object of each positive sample to create a negative one.

\todo{Ik denk dat we geen predicate corruption doen in \textsc{Synthology}}.

This baseline dataset is then used to train the RRN, which serves as the point of comparison for the model trained on \textsc{Synthology} data.

\subsection{Evaluation Metrics}

% do you intend to track the metrics with each training epoch, if so, mention this. Its important because comparing the metrics after the model has finished training on either dataset is only fair if it does the same amount of epochs on each training set

Because the RRN outputs a probability score $P := \P(\gen{s,p,o} = 1) \in [0, 1]$ indicating the likelihood of entailment, we evaluate the system as a \textit{binary classifier} rather than a ranking model (standard KGE metrics like Hits@K or MRR are thus not applicable). We focus on three key metrics.

Firstly, the Area Under the Curve (AUC) of the Receiver Operating Characteristic (ROC) is used to evaluate the model's overall threshold-independent capability to separate true inferences from false ones. Secondly, F1-scores are calculated to measure the harmonic mean of precision and recall at a standard threshold ($P > 0.5$). Finally, we specifically track the False Positive Rate (FPR) to evaluate the model's robustness against the `near-miss' logically false triples, which -- as discussed earlier -- are challenging for neuro-symbolic models.

% As we want to evaluate the training data rather than the model itself, does it make sense to add a metric like "training time/amount of epochs needed to reach certain F1 scores or FPR" or something along those lines?

\subsection{Hyperparameter Tuning \& Fair Comparison}
We vary customizable (YAML) configuration properties using automated Weights \& Biases (W\&B) \textit{sweeps}. A sweep is a systematic approach to hyperparameter optimization, where different configurations are tested and compared.

We ensure a fair comparison between the \textsc{Synthology}-generated dataset and the forward-chaining baseline, to make sure any differences in performance are due to the data generation \textit{method} rather than other factors. We provide both generators with the same entity pool size, the same total volume of input base facts ($|\mathcal{G}_\text{base}|$), and the same positive-to-negative target ratio $\rho_{\pm} = 1:1$.

\subsection{Open Source Implementation}
The entire codebase for \textsc{Synthology}, including the data generation engine, the training pipeline for the RRN, the ontology Turtle files, and the evaluation scripts, is made available as open source on \href{https://github.com/schependom/synthology}{GitHub}\footnote{\url{https://github.com/schependom/synthology}}. This allows anyone to reproduce our results, apply the generator to their own ontologies, and contribute improvements to the system.

\todo{Datasets should be available somewhere. Maybe mention seed?}

% =========================
% CHAPTER 5: EXPERIMENTS
% =========================
\section{Experiments \& Results}

Our experimental pipeline is designed to test three core hypotheses: the efficacy of our novel negative sampling, the structural superiority of backward-chained data for deep reasoning chains, and the tool's capacity to generalize across complex ontologies.

\subsection{Experiment 1: The Impact of Negative Sampling}
We first ablate the negative sampling strategies using the Family Tree ontology \cite{hohenecker2020ontology}. We construct three dataset variants using \textsc{Synthology}: pure random corruption, constrained random corruption, pure proof-based corruption. We train separate RRN models on each dataset and evaluate them on a skewed test set populated with logical near-misses. We anticipate that the model trained on proof-based negatives will exhibit a significantly lower False Positive Rate (FPR) when predicting complex relations, as the hard negatives prevent the network from adopting lazy topological heuristics.

\subsection{Experiment 2: Multi-Hop Reasoning Quality}
To show that our approach generates higher-quality multi-hop data, we isolate the RRN's performance on deep logic paths. We freeze a test set consisting exclusively of human-verified target facts requiring a proof depth of $\geq 3$ (e.g., \textit{hasGreatGrandparent}), as described in Appendix \ref{appendix:testset}. We then compare the performance of the RRN trained on a forward-chaining baseline dataset versus the RRN trained on \textsc{Synthology} data. We hypothesize that the \textsc{Synthology}-trained model will achieve a higher AUC-ROC because its training data guarantees complete logical traversal paths, whereas the forward-chaining nature of the baseline reasoner, based on the base facts, produces structurally shallow and disconnected graphs.

\subsection{Experiment 3: Generalization to Complex Ontologies}

To show that \textsc{Synthology} is ontology-agnostic and capable of scaling to highly complex, real-world schemas, we evaluate it against a modern, rigorously standardized benchmark. While historical benchmarks like the Lehigh University Benchmark (LUBM) are frequently cited, they contain existential quantifications that violate the strict OWL2 RL profile required for pure rule-based materialization. Therefore, we utilize OWL2Bench, a modern successor specifically engineered to test OWL2 RL compliant reasoners with advanced constructs such as complex property chains, inverse properties, and deep sub-class hierarchies.

\subsubsection{The Bottleneck of Standard Forward Chainers}

To establish a fair evaluation, \textsc{Synthology}'s backward-chaining approach must be compared against a ground-truth dataset generated by standard forward-chaining materialization. Initially, generic Semantic Web frameworks, such as Apache Jena, were considered for this baseline. However, empirical testing revealed a critical flaw in relying on standard rule engines for Machine Learning data preparation: they frequently fail to natively materialize complex multi-hop logic.

Specifically, when processing OWL2Bench, standard engines struggle to dynamically evaluate arbitrary \texttt{owl:propertyChainAxiom} constructs at runtime. Instead of yielding deep, multi-hop ABox targets (e.g., inferring a knows relationship through a 3-hop chain of courses and professors), the engine silently skips these chains and saturates the graph with trivial $\mathcal{O}(N^2)$ schema-level inequality assertions, such as \texttt{owl:differentFrom}. This highlights a severe infrastructure bottleneck in neuro-symbolic research: standard open-source tools cannot reliably generate the deep logical paths required to train models like the Recursive Reasoning Network (RRN).

\subsubsection{Establishing the Truth with Nemo}

To overcome this limitation and generate a mathematically perfect baseline, we utilize Nemo, a recently introduced toolkit for rule-based reasoning and data processing. Nemo's core is a highly scalable and efficient main-memory reasoner. Unlike legacy systems, Nemo natively supports an expressive extension of Datalog, including the handling of existential rules via a fast implementation of the restricted chase algorithm. Furthermore, supported formats and types for data values conform perfectly to the RDF standard, allowing it to ingest the OWL2Bench ontology and seamlessly materialize the exact multi-hop property chains that generic engines miss. Nemo's optimized storage can handle hundreds of millions of facts, making it the ideal state-of-the-art benchmark to represent the absolute limits of the forward-chaining paradigm.

\subsubsection{Experimental Design and Metrics}

In this experiment, we feed the raw OWL2Bench RL-profile TBox directly into \textsc{Synthology} without any manual rule configuration. We match the output yield of \textsc{Synthology} (Dataset B) to the exact number of target facts and Knowledge Graphs materialized by the Nemo baseline (Dataset A).

We then train separate instances of the RRN on both datasets and evaluate them on a frozen, held-out test set consisting exclusively of deep (3-hop and 4-hop) logical inferences. Beyond the model's predictive performance (measured via AUC-ROC and F1-Score), we track specific generation metrics to quantify \textsc{Synthology}'s efficiency.

\textit{Ontology Coverage} tracks the percentage of unique OWL2Bench rules successfully triggered during backward-chaining, proving the generator dynamically parses complex logic. \textit{Yield Rate} measures the ratio of base facts synthesized to inferred targets produced, measuring how efficiently \textsc{Synthology} constructs supporting evidence compared to the massive data dumps required by forward-chaining. Lastly, the generation time is recorded to demonstrate the practical feasibility of using \textsc{Synthology} for large-scale data synthesis.

By demonstrating that an RRN trained on \textsc{Synthology}'s engineered graphs outperforms one trained on Nemo's unguided forward-chained data -- while constrained to the exact same volume of targets -- we prove that \textsc{Synthology} provides a universally applicable, highly efficient solution for generating neuro-symbolic training data for any given OWL2 RL ontology.


% =========================
% CHAPTER 6: CONCLUSION
% =========================
\section{Discussion \& Future Work}
\notdone

\section{Conclusion}
\notdone

\newpage

\appendix
\section{Appendix: Test Set Construction for Experiment 2}
\label{appendix:testset}
\notdone

\newpage
\section{Appendix: \textsc{Synthology} Configuration Parameters}

\begin{table}[h!]
	\centering
	\begin{adjustbox}{rotate=90}
		\begin{tabular}{|l|l|l|l|p{9.76cm}|}
			\hline
			\textbf{YAML Parameter}            & \textbf{Symbol}      & \textbf{Type} & \textbf{Default} & \textbf{Description}                                                                                                                                                             \\
			\hline \hline
			\texttt{min\_individuals}          & $I_{\min}$           & int           & 1                & Lower acceptance bound on sample size: graphs with fewer individuals are rejected.                                                                                               \\
			\hline
			\texttt{max\_individuals}          & $I_{\max}$           & int           & 1000             & Upper acceptance bound on sample size: graphs with more individuals are rejected.                                                                                                \\
			\hline
			\texttt{min\_rules}                & $R_{\min}$           & int           & 1                & Minimum number of ontology rules selected per generated sample before proof generation.                                                                                          \\
			\hline
			\texttt{max\_rules}                & $R_{\max}$           & int           & 5                & Maximum number of ontology rules selected per generated sample.                                                                                                                  \\
			\hline
			\texttt{target\_min\_proofs\_rule} & $P_{\min}$           & int           & 5                & Target lower bound on proofs kept per selected rule; effectively bounded by how many valid proofs exist.                                                                         \\
			\hline
			\texttt{seed}                      & $s$                  & int           & 23               & Seed for pseudorandom sampling (rule selection, proof-root counts, and corruption choices), improving reproducibility.                                                           \\
			\hline

			\texttt{max\_recursion}            & $d_r$                & int           & 3                & Per-sample recursion cap for rule reuse in backward chaining; deeper recursion allows longer inference chains.                                                                   \\
			\hline
			\texttt{global\_max\_depth}        & $d_{\max}$           & int           & 10               & Absolute depth limit for recursive proof search; branches beyond this depth are pruned.                                                                                          \\
			\hline
			\texttt{max\_proofs\_per\_atom}    & $\kappa$             & int           & 5                & Hard cap on number of proofs emitted for one goal atom, preventing combinatorial explosion.                                                                                      \\
			\hline
			\texttt{individual\_pool\_size}    & $|\mathcal{U}|$      & int           & 60               & Target size of the reusable individual pool used when instantiating variables during proof construction.                                                                         \\
			\hline
			\texttt{individual\_reuse\_prob}   & $\pi_{\text{reuse}}$ & float         & 0.7              & Probability of reusing an existing individual from the pool rather than creating a new one.                                                                                      \\
			\hline
			\texttt{use\_signature\_sampling}  &                      & bool          & true             & If enabled, generated proofs are grouped by structural signature and one representative per group is sampled, improving diversity and reducing redundant Cartesian combinations. \\
			\hline
			\texttt{min\_proof\_roots}         & $U_{\min}$           & int           & 5                & Minimum number of independent root-generation cycles attempted per selected rule.                                                                                                \\
			\hline
			\texttt{max\_proof\_roots}         & $U_{\max}$           & int           & 20               & Maximum number of independent root-generation cycles attempted per selected rule.                                                                                                \\
			\hline
		\end{tabular}
	\end{adjustbox}
	\caption{Configuration parameters for ontology-based data generation in \textsc{Synthology} (1/2).}
	\label{tab:synthology-config}
\end{table}
\begin{table}[h!]
	\centering
	\begin{adjustbox}{rotate=90}
		\begin{tabular}{|l|l|l|l|p{9.76cm}|}
			\hline
			\textbf{YAML Parameter}         & \textbf{Symbol}     & \textbf{Type} & \textbf{Default} & \textbf{Description}                                                                                                                                     \\
			\hline
			\hline
			\texttt{always\_generate\_base} &                     & bool          & false            & If true, emits a base proof even when derivation rules apply; if false, base proofs are mainly used when no matching rule exists.                        \\
			\hline
			\texttt{min\_lcc\_ratio}        & $\rho_{\text{lcc}}$ & float         & 0.8              & Validation threshold for graph connectivity: the largest connected component must cover at least this fraction of individuals.                           \\
			\hline

			\texttt{strategy}               & $s_{-}$             & enum          & \texttt{mixed}   & Negative sampling mode: \texttt{random}, \texttt{constrained}, \texttt{type\_aware}, \texttt{proof\_based}, or \texttt{mixed}.                           \\
			\hline
			\texttt{ratio}                  & $\rho_{\pm}$        & float         & 1.0              & Target negative-to-positive ratio for generated examples; $\rho_{\pm}=1$ gives approximately balanced counts.                                            \\
			\hline
			\texttt{corrupt\_base\_facts}   &                     & bool          & false            & Enables corruption of proof-leaf base facts in proof-based logic; this controls whether propagated counterfactual negatives are produced in that branch. \\
			\hline
		\end{tabular}
	\end{adjustbox}
	\caption{Configuration parameters for ontology-based data generation in \textsc{Synthology} (2/2).}
	\label{tab:synthology-config}
\end{table}

% =========================
% BIBLIOGRAPHY
% =========================
\bibliographystyle{vancouver}

\newpage
\bibliography{references}


%%%%%%%%%%% The bibliography starts:
% \begin{thebibliography}{99}

% 	\bibitem{r1}

% \end{thebibliography}
\end{document}
